\section{Introduction}\label{sec:intro}

An index advisor chooses indexes for a workload by asking, for each candidate
configuration, what every query would cost if that configuration existed.
Relational systems answer with the optimizer's what-if interface, which costs a
query against hypothetical indexes that are never built~\cite{whatif1998}, and
advisors search over configurations with that interface as their
oracle~\cite{autoadmin1997,bruno2005relax,kossmann2020}. The oracle is an
estimate. Cardinality estimates in industrial optimizers are routinely and
substantially wrong~\cite{leis2015}, and optimizer estimates often predict that
an index will help a query that it in fact slows down~\cite{aimeetsai}. A recent
study on SQL Server found that index configurations recommended by large
language models, which often look worse than the industrial tuner's under the
optimizer's estimated cost, in several cases ran faster in actual execution,
although their quality varied widely~\cite{dtallm2026}. An advisor that ranks by
estimate can be outperformed by one that ignores the estimate.

Much analytical data now lives as open columnar files, typically Parquet, in
object storage, read by several engines rather than owned by
one~\cite{lakehouse2021}. Table formats record per-file statistics for data
skipping~\cite{deltalake,icebergspec}, and indexing subsystems add secondary
indexes for particular engines, in one case with what-if analysis over
hypothetical indexes inside that engine's optimizer~\cite{hyperspace}. The
storage tier itself has no optimizer, and so no what-if interface that every
engine reading the files can share, and none that is exact. Classical advisors
therefore have no engine-independent oracle there, and language-model tuners
have nothing trustworthy to be scored against.

We observe that this absence is an opportunity. A lake exposes its files. For
keyed workloads (look up a batch of keys, find the latest version of each key,
or find the version current at a timestamp), the I/O a run performs under a
hypothetical index depends on three things that can be recorded or read: the
run's key set, the row groups that contain each key, and the compressed size of
each row group. Replaying a recorded workload against row-group metadata
therefore computes, rather than estimates, the bytes, row groups, and files each
run would read under a candidate key specification (\adr{005}). We call this an
exact counterfactual what-if interface. Its exactness is scoped: it covers I/O
over files that exist, so a hypothetical layout is costed exactly only once it
is materialized, and costs above the scan come from measured execution
(Section~\ref{sec:exact}).

An exact oracle also changes what can be trusted from a language model.
LLM-based rewriters~\cite{genrewrite,llmr2,lithe,quite,e3rewrite,resequel}
generate candidates that must be scored and checked. Scoring by estimate
inherits the problem above. Checking with an equivalence
prover~\cite{cosette,qed,sqlsolver,verieql} covers restricted SQL fragments,
while pipeline steps are written in Spark, KQL, Scala, and configuration
(\adr{004}), and many useful changes are not query rewrites at all: an index, a
layout, a lookback bound. Bellwether therefore gives the model no authority. A
model may observe and propose; only the engine may measure, verify, and apply
(\adr{009}). The verifier is reconciliation of the pipeline's output over the
keys the workload touched (\adr{004}). This paper makes five contributions.

\begin{enumerate}[leftmargin=*,label=\textbf{C\arabic*}]
  \item \textbf{An exact counterfactual what-if interface for lakes.} Replay
    of recorded key sets against row-group metadata yields the exact I/O cost of
    hypothetical key specs (including composite and temporal keys) and plans
    over existing files; zero estimation error is a falsifiable design property
    (Section~\ref{sec:whatif}).
  \item \textbf{Reconciliation certificates.} Index-accelerated reconciliation
    by latest version, canonical row hash, and exact decimal totals yields
    certificates with an explicit coverage statement where SQL provers cannot
    apply; gating at the pipeline output admits non-local rewrites
    (Section~\ref{sec:cert}).
  \item \textbf{A hybrid proposer loop.} A reflective, population-based
    optimizer over plans, designs, and pipelines uses counterfactual attribution
    as its backward pass, a language model for structure, exact
    finite-difference descent for knobs, rule-driven enumeration, and Pareto
    selection across runs (Section~\ref{sec:evolve}).
  \item \textbf{A validated rulebook.} Distilled rules are promoted only when
    held-out replay reproduces a gain and the gate passes
    (Section~\ref{sec:rulebook}).
  \item \textbf{A trust boundary and a replayable ledger.} Every step is a
    content-addressed record, nothing a proposer emits is a score, and a search
    replays from its ledger byte for byte (Sections~\ref{sec:overview}
    and~\ref{sec:determinism}).
\end{enumerate}

This paper describes a fixed, documented design~\cite{bellwether} and a system
under construction; we cite design documents by file name and decision records
as ADR~001 to ADR~025. The contract, source, and ledger crates are merged, the
index core and two further crates are in review, and the optimizer is planned
(Section~\ref{sec:impl}). The evaluation methodology was fixed before any
implementation (\adr{019}); Section~\ref{sec:eval} reports it in full and marks
every result as pending. The paper's claims are therefore design properties,
each tied to a design document or an invariant with a named test, and none is a
measurement.
